\section{Notation}\label{sec:notation}

A consolidated reference for the mathematical objects used throughout this
report. Notation follows \citet{morimura2013} where it does, and is extended
in \cref{sec:methodology} where required.

\begin{table}[H]
\centering
\caption{Key symbols.}
\label{tab:notation}
\small
\begin{tabular}{@{}lp{0.74\linewidth}@{}}
\toprule
\textbf{Symbol} & \textbf{Meaning} \\
\midrule
$\mathcal{X}$, $n$ & Full state space (directed road edges); $n = |\mathcal{X}|$. \\
$\Xo \subset \mathcal{X}$ & Subset of \emph{observed} states at which counts are available. \\
$\mathrm{adj}(x)$ & Set of legal successors of state $x$ (out-neighbours in the directed graph). \\
$\pT(x'\mid x)$ & One-step transition probability from $x$ to $x'$; row $x$ of the kernel $\Pchain$. \\
$\pI(x)$ & Initial / teleport probability on $\mathcal{X}$; row vector. \\
$\beta \in [0,1]$ & Continuation probability per step; together with $\pI$ defines the stationary mixture $P = \beta\Pchain + (1-\beta)\,\bm{1}\pI^\top$. \\
$\pi$ & Stationary distribution of $P$; obtained from $\pi^\top P = \pi^\top$, $\bm{1}^\top\pi = 1$. \\
$\fobs(x)$ & Empirical visit count (or normalised frequency) at $x$. Observed only for $x \in \Xo$. \\
$\gobs(x, x')$ & Empirical (or analytical) \emph{hitting rate} from $x$ to $x'$, a discounted first-passage probability on $\Xo \times \Xo$. \\
$\phi_T(x') \in \mathbb{R}^{d_T}$ & Destination-state feature vector (\cref{sec:features}). \\
$\psi(x, x') \in \mathbb{R}^{d_\psi}$ & Directed-edge feature vector. \\
$\nu \in \mathbb{R}^n$ & Per-state base parameters of $\pI$. \\
$\omega^{\mathrm{loc}} \in \mathbb{R}^E$ & Per-edge local weights of $\pT$, with $E = \sum_x |\mathrm{adj}(x)|$ the number of directed edges. \\
$\omega^{\mathrm{glo}_1} \in \mathbb{R}^{d_T}$ & Shared weights on destination features $\phi_T$ (paper Eq.~17). \\
$\omega^{\mathrm{glo}_2} \in \mathbb{R}^{d_\psi}$ & Shared weights on edge features $\psi$ (paper Eq.~17). \\
$\theta$ & Full parameter vector $[\nu;\omega^{\mathrm{loc}};\omega^{\mathrm{glo}_1};\omega^{\mathrm{glo}_2}] \in \mathbb{R}^{d}$, $d = n + E + d_T + d_\psi$. \\
$\lambda$ & Tikhonov regularisation strength on $\|\theta\|^2$. \\
$\gamma \in [0,1]$ & Mixing weight between the stationary-match loss $L_d$ and the hitting-rate loss $L_h$. \\
$\Lam(\tau)$ & Log-likelihood-ratio score of a test trajectory $\tau$ under two fitted chains (\cref{sec:methodology}). \\
$\AUC$ & Area under the ROC curve of the LR score against labelled regime test trajectories. \\
$\RMAE(\hat{\pi},\pi)$ & Relative mean absolute error of recovered stationary mass at unobserved states. \\
\bottomrule
\end{tabular}
\end{table}

Throughout, ``regime'' refers to one of two distinguishable driver
populations whose chains we fit and compare. In controlled simulation
this is intrinsic vs.\ sat-navigation-influenced; in real-world data it
is a behavioural proxy (e.g.\ rush-hour vs.\ off-peak). The contrast
$\Delta_{\mathrm{fit}}(x, \cdot) = \pT^{(B)}(x,\cdot) - \pT^{(A)}(x,\cdot)$
between two fitted chains is the central quantity assessed against the
empirical ground-truth contrast
$\Delta_{\mathrm{emp}}(x, \cdot)$.
